\begin{flushleft}
	A detecção de objetos em tempo real é fundamental para o avanço de sistemas autônomos e de segurança. A integração dessa tecnologia a gimbals multi-eixos não apenas proporciona estabilização de imagem e rastreamento preciso, mas também expande significativamente o campo de visão operacional. No contexto aeronáutico, essa integração se manifesta em sensores optoeletrônicos cruciais para operações militares e de reconhecimento. Contudo, um desafio persistente reside na manutenção do rastreamento confiável, dado o intrínseco equilíbrio entre a precisão do algoritmo de detecção e o custo computacional, bem como as dinâmicas complexas do gimbal (incluindo não linearidades e atrasos de comunicação).
	Este projeto de mestrado tem como objetivo principal desenvolver uma solução de controle robusta para aplicações de rastreamento de objetos em imagens em tempo real. Para tanto, emprega-se o YOLOv8, um algoritmo de detecção de última geração, cuja saída alimenta um sistema de controle PID. A atuação do controle é sobre um gimbal de dois graus de liberdade, responsável por ajustar a orientação de uma câmera para manter o objeto centralizado. A validação do desempenho das técnicas de controle propostas é realizada em um ambiente de simulação (Simulink), executado em tempo real, permitindo uma avaliação abrangente sob condições que emulam cenários operacionais.
\end{flushleft}